% =============================================================================
% SVVVL Research Paper
% An Evolutionary Foundation for Continuous-Time Asset Allocation
% =============================================================================

\documentclass[11pt]{article}

% ----- Page geometry -------------------------------------------------------
\usepackage[
  letterpaper,
  top=0.9in,
  bottom=0.9in,
  left=0.85in,
  right=0.85in,
  headheight=14pt
]{geometry}

% ----- Encoding & language -------------------------------------------------
\usepackage[utf8]{inputenc}
\usepackage[T1]{fontenc}
\usepackage[english]{babel}



% ----- Typography ----------------------------------------------------------
\usepackage{XCharter}

% Section headings
\usepackage{titlesec}
\titleformat{\section}
  {\normalfont\Large\bfseries}{\thesection}{0.6em}{}
\titleformat{\subsection}
  {\normalfont\large\bfseries}{\thesubsection}{0.6em}{}
\titlespacing*{\section}{0pt}{1.4em}{0.6em}
\titlespacing*{\subsection}{0pt}{1.1em}{0.4em}

% Remove section numbering
\setcounter{secnumdepth}{0}

% ----- Spacing & paragraphs ------------------------------------------------
\usepackage{setspace}
\setstretch{1.15}
\setlength{\parskip}{0.6em}
\setlength{\parindent}{0pt}

% ----- Graphics & color ----------------------------------------------------
\usepackage{graphicx}
\usepackage{xcolor}
\definecolor{linkcolor}{HTML}{1F4E79}

% ----- Hyperlinks ----------------------------------------------------------
\usepackage[
  colorlinks=true,
  linkcolor=linkcolor,
  urlcolor=linkcolor,
  citecolor=linkcolor
]{hyperref}

% ----- Tables, figures, captions ------------------------------------------
\usepackage{booktabs}
\usepackage{array}
\usepackage{caption}
\captionsetup{font=small, labelfont=bf, skip=6pt}

% ----- Math ----------------------------------------------------------------
\usepackage{amsmath, amssymb}
\usepackage{amsthm}
\usepackage{mathtools}

% ----- Footnotes -----------------------------------------------------------
\usepackage[hang, flushmargin]{footmisc}
\setlength{\footnotesep}{0.8em}

% ----- Headers & footers ---------------------------------------------------
\usepackage{fancyhdr}
\pagestyle{fancy}
\fancyhf{}
\renewcommand{\headrulewidth}{0pt}
\fancyfoot[C]{\thepage}

\fancypagestyle{firstpage}{%
  \fancyhf{}%
  \renewcommand{\headrulewidth}{0pt}%
}


% ----- Logo lockup ---------------------------------------------------------
\newcommand{\WORDMARKFILE}{../assets/SVVVL_Research_Wordmark.pdf}
\newcommand{\WORDMARKWIDTH}{1.5in}
\newcommand{\LOCKUPGAP}{0.01in}

% The lockup itself: symbol above wordmark, both centered within a box
% the width of the symbol so the wordmark sits naturally beneath it.
\newcommand{\logolockup}{%
  \begin{minipage}{\WORDMARKWIDTH}
    \centering
    \includegraphics[width=\WORDMARKWIDTH]{\WORDMARKFILE}
  \end{minipage}%
}

% Title block command: \makecover{Title}{Subtitle}{Authors}{Date}
\newcommand{\makecover}[4]{%
  \thispagestyle{firstpage}%
  \noindent\logolockup\par
  \vspace{2.2em}
  {\fontsize{26pt}{30pt}\selectfont\bfseries
   \hyphenpenalty=10000\exhyphenpenalty=10000\sloppy
   #1\par}
  \vspace{0.6em}
  {\fontsize{15pt}{18pt}\selectfont\itshape\color{black!70}
   \hyphenpenalty=10000\exhyphenpenalty=10000\sloppy
   #2\par}
  \vspace{1.6em}
  {\normalsize\color{black!60} #3\par}
  \vspace{0.2em}
  {\normalsize\color{black!60} #4\par}
  \vspace{1.8em}
}

% ----- Theorem environments ------------------------------------------------
\newtheorem{proposition}{Proposition}
\newtheorem{remark}{Remark}
\newtheorem{conjecture}{Conjecture}

% ----- Custom commands -----------------------------------------------------
\newcommand{\lopt}{\ell^{*}}
\newcommand{\lkelly}{\ell^{**}}
\newcommand{\mue}{\mu_{e}}
\newcommand{\mub}{\mu_{b}}
\newcommand{\mus}{\mu_{s}}
\newcommand{\E}{\mathbb{E}}
\newcommand{\Var}{\mathrm{Var}}
\newcommand{\Corr}{\mathrm{Corr}}
\newcommand{\dd}{\mathrm{d}}

% =============================================================================
\begin{document}
% =============================================================================

% --- COVER -------------------------------------------------------------------
\makecover
  {Nature Abhors an Undiversified Bet}
  {An Evolutionary Foundation for Continuous-Time Asset Allocation}
  {J.~E. Whitehead~III\footnote{SVVVL, Towson, MD, United States.
   Contact: \href{mailto:james.whitehead@svvvl.com}{james.whitehead@svvvl.com}.
   See Disclosures at the end of this document.}}
  {\today}

% --- EPIGRAPH ----------------------------------------------------------------
\vspace{-0.5em}
\begin{quote}
\itshape
``Nature abhors an undiversified bet.''\\
\textnormal{--- Brennan and Lo (2012)}
\end{quote}
\vspace{0.5em}

% --- ABSTRACT ----------------------------------------------------------------
\begin{quote}
\itshape
We extend the evolutionary binary-choice framework of Brennan and
Lo~[\hyperlink{ref:3}{3}] to continuous-time asset allocation with one risky and one
risk-free asset. The key move is a change in the unit of selection:
individual dollars, not investors. Every dollar faces the same
systematic Brownian shock, so the Brennan-Lo logic applies
directly: the growth-optimal strategy is randomization. Maximizing
instantaneous log-portfolio growth yields
$\lopt = \mue/\sigma^{2} + (1/\sigma)\,\dd W_{t}/\dd t$.
The deterministic part is the Kelly criterion~[\hyperlink{ref:7}{7}] and, in the
continuous-time limit, Merton's rule~[\hyperlink{ref:9}{9},\hyperlink{ref:10}{10}]. The stochastic part is the investor's belief about market
randomness, characterized by its correlation $\rho$ with the true
Brownian innovation, the continuous-time analog of the Brennan-Lo
intelligence parameter. Substituting $\lopt$ into the log-portfolio
dynamics produces an intelligence correction of $\rho -
\tfrac{1}{2} \in [-\tfrac{3}{2}, +\tfrac{1}{2}]$: misaligned beliefs
destroy three times what aligned beliefs create. The stochastic term
also resolves a long-standing puzzle. The instantaneous
growth-optimal rule $\lopt$ and the long-run rule $\lkelly =
\mue/\sigma^{2}$ are distinct objects. The variance of the
randomization term, $1/(\dd t \cdot \sigma^{2})$, is the hedge. The
expected allocation remains Kelly. The spread around it widens as the
horizon shortens. Because no investor lives forever, that spread is
always positive.
\end{quote}


\newpage
% --- INTRODUCTION ------------------------------------------------------------
\section{Introduction}

The Kelly criterion~[\hyperlink{ref:7}{7}] and Merton's portfolio rules~[\hyperlink{ref:9}{9},\hyperlink{ref:10}{10}] are
derived from optimization principles that, while mathematically
clean, leave open the question of why evolution would produce
investors who behave consistently with them. They tell us what such
investors would do, not why selection should favor them.

Brennan and Lo~[\hyperlink{ref:3}{3}] asked a harder question: what does natural
selection actually produce? In their binary-choice model, a population
of organisms picks between two actions. Each action yields a random
number of offspring. The punchline is quiet but consequential. When
environmental shocks are systematic, when the whole population is in the same ecological boat, rational deterministic behavior loses.
Randomization wins. Probability matching, loss aversion, and other
phenomena that look irrational from the individual's perspective are
optimal from the population's. Rational behavior emerges only when
shocks are idiosyncratic.

This paper carries that logic into continuous-time portfolio
allocation. One risky asset, one risk-free asset. The risky asset
exposes every investor to the same Brownian shock simultaneously.
That is complete systematic correlation, and it appears at first to
rule out any deterministic portfolio rule. A single reframing
dissolves the difficulty.

We make four contributions.

\textbf{Contribution 1: The population-of-dollars reframing.}
At first glance, the Brennan-Lo logic seems to apply to investors.
The original framework describes creatures that face a binary choice,
and investors are the natural analog. But squint, tilt the head, and
the logic admits a different reading. It applies just as cleanly to
the dollars inside an investor's portfolio. The investors are no
longer center stage; their dollars are. Dollars are the population
for the investor, just as investors are the population for the
market. With this reframing, each dollar faces the same systematic
Brownian shock $\dd W_{t}$, the Brennan-Lo logic for systematic
environments applies directly, and the growth-optimal strategy is
randomization. Each dollar draws its allocation from a Gaussian
distribution, a biased coin flip, and the aggregate Kelly weight
emerges from the law of large numbers.

\textbf{Contribution 2: Kelly and Merton from selection.}
Maximizing instantaneous log-portfolio growth yields
\[
  \lopt \;=\; \frac{\mue}{\sigma^{2}} + \frac{1}{\sigma}
  \frac{\dd W_{t}}{\dd t},
\]
where $\mue = \mus - \mub$ is the expected excess return. The deterministic
part is the Kelly weight, which is Merton's fraction in the continuous-time
limit. Both results fall out of evolutionary selection.

\textbf{Contribution 3: Intelligence, defined.}
Write the investor's belief about market randomness as $M_{t}'$,
decomposed as
\[
  M_{t}' \;=\; \sqrt{1-\rho^{2}}\,\widetilde{W}_{t}' + \rho\, W_{t}',
\]
where $\rho = \Corr[M_{t}', W_{t}']$ measures how well the
investor's model of market randomness tracks the truth. Substitute
$\lopt$ into the log-portfolio dynamics. The drift picks up an
intelligence correction of $\rho - \tfrac{1}{2} \in
[-\tfrac{3}{2}, +\tfrac{1}{2}]$. Bad beliefs destroy three times
what good beliefs create.

\textbf{Contribution 4: The Kelly short-run puzzle, resolved.}
$\lopt = \mue/\sigma^{2} + (1/\sigma)\,\dd W_{t}/\dd t$ and $\lkelly = \mue/\sigma^{2}$ are different objects.
The first maximizes instantaneous stochastic growth rate of the portfolio. The second
maximizes long-run average growth rate. The stochastic term is a hedge:
the expected dollar allocation is always Kelly, but each dollar draws
from a Gaussian centered at Kelly with variance $1/(\dd t\cdot\sigma^{2})$.
The shorter the horizon, the wider the spread, and the greater the
hedge required. The pure Kelly weight $\lkelly$ is the limiting case
as $\dd t \to \infty$, a horizon no investor reaches.

The paper proceeds as follows. Section~2 surveys the literature.
Section~3 sets up the model. Section~4 derives the optimal allocation.
Section~5 develops the population-of-dollars reframing. Section~6
maps the framework to Brennan-Lo. Section~7 defines intelligence.
Section~8 treats bounded rationality. Section~9 resolves the Kelly
short-run puzzle. Section~10 discusses implications. Section~11
concludes.

% --- RELATED LITERATURE ------------------------------------------------------
\section{Related Literature}

\subsection{Evolutionary foundations of economic behavior}

Brennan and Lo~[\hyperlink{ref:3}{3}] build a binary-choice model in which a population
of organisms picks between two actions, each yielding a random number
of offspring. The growth-optimal behavior $f^{*}$ maximizes the
geometric growth rate $\alpha(f) = \E[\log(f x_{a} + (1-f)x_{b})]$.
The central result is a clean dichotomy: systematic shocks favor
randomization; idiosyncratic shocks favor individually-optimal
deterministic behavior. Intelligence is $\rho = \Corr[I^{f}_{it},
(x_{a} - x_{b})]$: how well the organism's choice tracks which
action nature rewards. When intelligence carries a cost $c(\rho)$,
bounded rationality is the outcome. The earlier model is in Brennan
and Lo~[\hyperlink{ref:2}{2}]. The adaptive coin-flipping interpretation of random
individual variation under natural selection originates with Cooper
and Kaplan~[\hyperlink{ref:4}{4}], whose decision-theoretic framework provides the
conceptual foundation for the randomization results developed here.
Related work on the natural selection of utility
functions~[\hyperlink{ref:12}{12},\hyperlink{ref:13}{13},\hyperlink{ref:14}{14},\hyperlink{ref:15}{15}] and the evolution of trading
strategies~[\hyperlink{ref:5}{5}] informs the broader context.

\subsection{The Kelly criterion and growth-optimal portfolios}

Kelly~[\hyperlink{ref:7}{7}] found the growth-optimal betting strategy for an investor
with an edge. Breiman~[\hyperlink{ref:1}{1}] proved it wins almost surely in the long
run. Thorp~[\hyperlink{ref:16}{16}] and Cover~[\hyperlink{ref:6}{6}] extended the result to portfolios.
Every paper in this line assumes growth-optimality as the goal. We
derive it.

\subsection{Merton's continuous-time framework}

Merton~[\hyperlink{ref:9}{9},\hyperlink{ref:10}{10}] solved the dynamic portfolio problem in continuous
time under general utility. Under log utility, the solution collapses
to a clean myopic rule identical to Kelly in the continuous-time
limit. We supply the evolutionary argument for why the market
produces investors whose behavior is consistent with this solution.

\subsection{Ergodicity economics}

Peters~[\hyperlink{ref:11}{11}] makes the case that time-average growth, not
ensemble-average growth, is the right objective. The time-average
growth rate of a leveraged portfolio is
\[
  \bar{g}(\ell) \;=\; \mub + \ell\mue - \frac{\ell^{2}\sigma^{2}}{2},
\]
with optimal leverage $\ell_{\mathrm{opt}} = \mue/\sigma^{2}$,
derived from the time-average growth rate directly. We embed this
result in the Brennan-Lo evolutionary
framework. The variance drag $\ell^{2}\sigma^{2}/2$ is the
continuous-time face of the same evolutionary penalty that makes
over-commitment fatal in Brennan-Lo.

\subsection{The Adaptive Markets Hypothesis}

Lo~[\hyperlink{ref:8}{8}] reads financial markets as adaptive systems in which
strategies compete and evolve. The intelligence parameter $\rho$ and
its cost $c(\rho)$ developed here give that intuition a quantitative
backbone: market efficiency is time-varying because the cost-benefit
calculus of belief accuracy shifts with market conditions.

% --- THE MODEL ---------------------------------------------------------------
\section{The Continuous-Time Portfolio Model}

\subsection{Asset dynamics}

Two assets. The risky asset $S$ follows a geometric Brownian motion,
\begin{equation}
  \frac{\dd S}{S} \;=\; \mus\,\dd t + \sigma\,\dd W_{t},
  \label{eq:risky}
\end{equation}
where $\mus > 0$ is the drift, $\sigma > 0$ is the volatility, and
$W_{t}$ is a standard Brownian motion\footnote{Throughout the paper,
we adopt continuous-time notation for derivational convenience. The
formal expression $\dd W_{t}/\dd t$ that appears later in the
derivation is not a classical derivative; $W_{t}$ is nowhere
differentiable in the ordinary sense. We use it as shorthand for the
finite increment $\Delta W_{t}/\Delta t \sim \mathcal{N}(0,
1/\Delta t)$ at a small but strictly positive rebalancing interval
$\Delta t > 0$. The continuous-time framing serves as the idealized
limit useful for derivation; the result is meant to be applied at
finite $\Delta t$. This is consistent with the long-standing
practice in physics, control theory, and quantitative finance of
treating white noise as the formal derivative of a Wiener process,
with all expressions interpreted at a finite sampling frequency
when implemented. It is also consistent with the argument developed
in Section~9: $\dd t \to \infty$ is never reached in practice, and
neither is $\dd t \to 0$.} on a filtered probability space
$(\Omega, \mathcal{F}, \{\mathcal{F}_{t}\}, \mathbb{P})$. The
risk-free asset $B$ grows at a known rate,
\begin{equation}
  \frac{\dd B}{B} \;=\; \mub\,\dd t.
  \label{eq:riskfree}
\end{equation}
The excess drift is
\begin{equation}
  \mue \;=\; \mus - \mub \;>\; 0,
  \label{eq:excess}
\end{equation}
and we assume the risky asset earns a positive premium throughout.

\subsection{Portfolio dynamics}

Let $\ell \in \mathbb{R}$ be the portfolio weight on the risky asset.
We allow $\ell$ to range over the real line. Short positions
($\ell < 0$), full investment ($\ell = 1$), and leverage ($\ell > 1$)
are all permitted. This generality is required: the optimal allocation
$\lopt$ derived in Section~4 is drawn from a distribution supported
on $\mathbb{R}$. The portfolio return is
\begin{equation}
  \frac{\dd P}{P} \;=\; \ell\,\frac{\dd S}{S}
  + (1-\ell)\,\frac{\dd B}{B}.
  \label{eq:portfolio}
\end{equation}
Substituting \eqref{eq:risky} and \eqref{eq:riskfree},
\begin{equation}
  \dd P \;=\; P\bigl(\mub + \mue\,\ell\bigr)\dd t
  + P\,\ell\,\sigma\,\dd W_{t}.
  \label{eq:portfoliosde}
\end{equation}
It\^{o}'s lemma applied to $\ln P(t)$ gives
\begin{equation}
  \dd\ln P(t) \;=\; \Bigl(\mub + \mue\,\ell
  - \frac{\sigma^{2}}{2}\,\ell^{2}\Bigr)\dd t
  + \sigma\,\ell\,\dd W_{t}.
  \label{eq:logportfolio}
\end{equation}
The term $\sigma^{2}\ell^{2}/2$ is the It\^{o} variance drag. It is
the continuous-time price of holding the risky asset, and it cannot
be avoided. Integrating \eqref{eq:logportfolio},
\begin{equation}
  \ln P(t) \;=\; \ln P(0) + \Bigl(\mub + \mue\,\ell
  - \frac{\sigma^{2}}{2}\,\ell^{2}\Bigr)t
  + \sigma\,\ell\,W_{t},
  \label{eq:logwealth}
\end{equation}
so the time-average growth rate is
\begin{equation}
  \bar{g}(\ell) \;=\; \mub + \mue\,\ell
  - \frac{\sigma^{2}}{2}\,\ell^{2}.
  \label{eq:growthrate}
\end{equation}
A concave quadratic in $\ell$. Peters~[\hyperlink{ref:11}{11}] arrives at the same
expression from ergodicity arguments.

% --- OPTIMAL ALLOCATION ------------------------------------------------------
\section{The Optimal Allocation Rule}

\subsection{Maximizing instantaneous log-portfolio growth}

Differentiate $\dd\ln P(t)$ in \eqref{eq:logportfolio} with respect
to $\ell$ and set to zero,
\begin{equation}
  \frac{\partial\,\dd\ln P}{\partial\,\ell}
  \;=\; \mue\,\dd t - \sigma^{2}\,\ell\,\dd t
  + \sigma\,\dd W_{t} \;=\; 0.
  \label{eq:foc}
\end{equation}
Solving for $\ell$,
\begin{equation}
  \boxed{
  \lopt \;=\; \frac{\mue}{\sigma^{2}}
  + \frac{1}{\sigma}\,\frac{\dd W_{t}}{\dd t}.
  }
  \label{eq:lstar}
\end{equation}
Two terms. The first, $\mue/\sigma^{2}$, is the Kelly weight: excess return over variance. The second,
$(1/\sigma)\,\dd W_{t}/\dd t$, is stochastic. The optimal allocation
moves with the Brownian innovation at every instant. It is a process,
not a number.

\subsection{The belief decomposition}

The stochastic term is the investor's response to market randomness.
Call it her belief. Define
\begin{equation}
  M_{t}' \;\equiv\; \frac{\dd W_{t}}{\dd t},
  \label{eq:Mprime}
\end{equation}
the investor's view of the market's instantaneous innovation. A
perfectly informed investor sets $M_{t}' = W_{t}'$, where
$W_{t}' \equiv \dd W_{t}/\dd t$. Most investors do not. We write
\begin{equation}
  M_{t}' \;=\; \sqrt{1-\rho^{2}}\,\widetilde{W}_{t}'
  + \rho\,W_{t}',
  \label{eq:belief}
\end{equation}
where $\widetilde{W}_{t}' \sim \mathcal{N}(0, 1/\dd t)$ may be regarded as
the investor's view on the non-systematic part of the risky asset's return
$W_{t}' \sim \mathcal{N}(0, 1/\dd t)$. The correlation
\begin{equation}
  \rho \;=\; \Corr[M_{t}', W_{t}'] \;\in\; [-1, 1]
  \label{eq:rho}
\end{equation}
measures how well the investor reads the market. The optimal
allocation becomes
\begin{equation}
  \lopt \;=\; \frac{\mue}{\sigma^{2}}
  + \frac{1}{\sigma}\,M_{t}'.
  \label{eq:lstarM}
\end{equation}

% --- POPULATION OF DOLLARS ---------------------------------------------------
\section{The Population of Dollars}

\subsection{Why investors are the wrong population}

With one risky asset, every investor sees the same $\dd W_{t}$ at
the same time. That is complete systematic correlation. The
Brennan-Lo result~[\hyperlink{ref:3}{3}] is unambiguous: under systematic shocks, no
fixed deterministic allocation survives. A bad enough realization of
$\dd W_{t}$ will eventually ruin any investor holding a fixed $\ell$.
The answer is to find the right population, not to escape the
systematic risk.

\subsection{The Gaussian coin flip}

The population is individual dollars, not investors. Each dollar is
the unit of selection. Every dollar faces the same systematic
$\dd W_{t}$. The Brennan-Lo logic applies to the dollar population
directly, and the conclusion is the same as it always is under
systematic risk: randomize. Each dollar $i$ draws its allocation
independently,
\begin{equation}
  \lopt_{i} \;\sim\; \mathcal{N}\!\left(
    \frac{\mue}{\sigma^{2}},\;
    \frac{1}{\dd t\cdot\sigma^{2}}
  \right).
  \label{eq:gaussiancoin}
\end{equation}
The mean is the Kelly weight $\mue/\sigma^{2}$;
the variance is $1/(\dd t\cdot\sigma^{2})$, the Gaussian coin flip.
Probability matching in Brennan-Lo and this coin flip are the same
thing, expressed in different notation. Both are instances of the
adaptive coin-flipping principle of Cooper and Kaplan~[\hyperlink{ref:4}{4}].

\subsection{The aggregate result}

The draws are independent across dollars. The law of large numbers
takes over,
\begin{equation}
  \E[\lopt_{i}] \;=\; \frac{\mue}{\sigma^{2}} \;=\; \lkelly.
  \label{eq:aggregate}
\end{equation}
The Kelly weight appears at the portfolio level as the aggregate of
optimal dollar-level randomization. In Brennan-Lo, probability
matching at the individual level produces the fastest-growing
population. Here, Gaussian randomization at the dollar level produces
the Kelly portfolio. Same logic, continuous time.

% --- MAPPING -----------------------------------------------------------------
\section{Mapping to the Brennan-Lo Framework}

The correspondence between the continuous-time framework and the
Brennan-Lo binary-choice model is exact. Table~\ref{tab:mapping}
lays it out.

\begin{table}[ht]
\centering
\caption{Mapping between the Brennan-Lo framework and the
continuous-time portfolio model.}
\label{tab:mapping}
\renewcommand{\arraystretch}{1.4}
\begin{tabular}{p{6.5cm} p{6.5cm}}
\toprule
\textbf{Brennan-Lo (2012)} & \textbf{This Paper} \\
\midrule
Individual organism & Individual dollar \\
Population of organisms & Portfolio of dollars \\
Behavioral phenotype $f \in [0,1]$ & Portfolio weight $\ell \in \mathbb{R}$ \\
Offspring count $x_{a}$, $x_{b}$ & Gross portfolio return \\
Geometric growth rate $\alpha(f)$ & Log-portfolio growth $\dd\ln P(t)$ \\
Systematic shock (single niche) & Shared Brownian increment $\dd W_{t}$ \\
Randomization under systematic risk & Gaussian coin flip \eqref{eq:gaussiancoin} \\
Growth-optimal phenotype $f^{*}$ & Kelly weight $\lkelly = \mue/\sigma^{2}$ \\
Intelligence $\rho = \Corr[I^{f}_{it}, (x_{a}-x_{b})]$ & $\rho = \Corr[M_{t}', W_{t}']$ \\
Volatility drag (implicit) & It\^{o} drag $\sigma^{2}\ell^{2}/2$ \\
\bottomrule
\end{tabular}
\end{table}

The key correspondence: $\alpha(f)$ in Brennan-Lo and $\dd\ln P(t)$
here are both the objects selection maximizes. The It\^{o} variance
drag $\sigma^{2}\ell^{2}/2$ is the continuous-time face of the same
penalty that makes over-commitment lethal in Brennan-Lo. Nature
penalizes the undiversified bet. In Brennan-Lo the penalty is
extinction. Here it is the volatility drag, smooth and analytic, but no
less real.

% --- INTELLIGENCE ------------------------------------------------------------
\section{Intelligence: A Continuous-Time Definition}

\subsection{The Brennan-Lo definition}

Brennan and Lo~[\hyperlink{ref:3}{3}] define intelligence as the correlation between an
organism's choice and the difference in reproductive outcomes,
\[
  \rho \;=\; \Corr\!\bigl[I^{f}_{it},\; (x_{a} - x_{b})\bigr].
\]
Populations with $\rho > 0$ grow faster. Populations with $\rho < 0$
do not survive. Intelligence, in this framework, is alignment between
behavior and what nature rewards.

\subsection{The continuous-time analog}

Here, intelligence is the correlation between the investor's belief
process and the true market innovation,
\[
  \rho \;=\; \Corr[M_{t}', W_{t}'].
\]
An investor is intelligent to the degree that her model of market
randomness tracks the truth. The definition carries the precise
evolutionary meaning of Brennan-Lo, behavior correlated with good outcomes, expressed in continuous time.

\subsection{The optimal log-portfolio dynamics under $\rho$}

Substitute $\lopt$ from \eqref{eq:lstarM} into \eqref{eq:logportfolio}
using the decomposition \eqref{eq:belief},
\begin{equation}
  \dd\ln P(t) \;=\; \left(\mub
  + \frac{1}{2}\left(\frac{\mue}{\sigma}\right)^{\!2}\right)\dd t
  + \left(\rho - \frac{1}{2}\right)
  + \frac{\mue}{\sigma}\,\dd W_{t}.
  \label{eq:optlogdyn}
\end{equation}
Three terms. The first scales with time: the risk-free rate and the
squared Sharpe ratio. The second is the intelligence correction
$\rho - \tfrac{1}{2}$. The third is the diffusion. Algebraically,
the correction enters as a level shift outside $\dd t$. We will
return to the question of whether $\rho$ itself can be treated as
independent of $\dd t$ shortly; physically, it cannot.

\subsection{A conjecture: $\rho$ as a function of the clock}

Treating $\rho$ as a free parameter on $[-1, 1]$, independent of
$\dd t$, is mathematically clean but physically suspect. Any real
strategy that converts information into edge, into positive $\rho$, faces hard limits. There is only so much information any
agent can process per unit time, and only so much of that information
that translates into alignment with $W_{t}'$. As the horizon $\dd t$
shrinks, the maximum achievable $\rho$ should shrink with it. The
shorter the window, the less time there is to gather, process, and
act on signal.

We therefore conjecture an asymmetric, horizon-dependent bound on
$\rho$.

\begin{conjecture}[Horizon-dependent bounds on intelligence]\label{conj:rhomax}
The achievable correlation $\rho$ between the investor's belief
process $M_{t}'$ and the true innovation $W_{t}'$ is bounded by
\[
  \rho_{\min} \;\leq\; \rho \;\leq\; \rho_{\max}(\dd t),
\]
where
\[
  \rho_{\min} \;=\; -1 \quad\text{(constant)},
  \qquad
  \rho_{\max}(\dd t) \;\in\; [0,1]
  \;\text{is monotonically increasing in } \dd t.
\]
\end{conjecture}

The asymmetry is the point. Building positive $\rho$ requires
information, processing, and execution, all rate-limited. Negative
$\rho$ requires none of these. An investor can achieve perfect
misalignment instantaneously, by holding beliefs that are exactly
backward. Stupid knows no speed limit.

The conjecture sharpens \eqref{eq:optlogdyn} into a horizon-dependent
statement. Substituting $\rho_{\max}(\dd t)$ for the upper end of
the range, the achievable intelligence correction runs from
$-\tfrac{3}{2}$ up to $\rho_{\max}(\dd t) - \tfrac{1}{2}$, which
itself shrinks toward $-\tfrac{1}{2}$ as $\dd t \to 0$. At short
horizons, the upside from intelligence narrows; the downside does
not. We treat \eqref{eq:optlogdyn} as the reduced-form expression
and Conjecture~\ref{conj:rhomax} as the structural constraint that
must eventually be respected when implementing the framework against
real markets.

\subsection{The full range of $\rho \in [-1, 1]$}

Setting the conjecture aside for the moment and reading $\rho$
purely as a parameter, the intelligence correction $\rho -
\tfrac{1}{2}$ runs from $-\tfrac{3}{2}$ to $+\tfrac{1}{2}$.
Table~\ref{tab:rho} gives the full picture. The reader should keep
Conjecture~\ref{conj:rhomax} in mind: the upper end of the table is
not unconditionally available at every horizon.

\begin{table}[ht]
\centering
\caption{Intelligence correction $\rho - \tfrac{1}{2}$ for
$\rho \in [-1, 1]$.}
\label{tab:rho}
\renewcommand{\arraystretch}{1.4}
\begin{tabular}{ccc p{7cm}}
\toprule
$\rho$ & $M_{t}'$ & Correction & Interpretation \\
\midrule
$1$ & $W_{t}'$ & $+\tfrac{1}{2}$ &
  Perfect alignment; maximum growth premium \\
$\left(\tfrac{1}{2}, 1\right)$ & Partial signal & $\left(0, +\tfrac{1}{2}\right)$ &
  Bounded intelligence; partial signal and noise \\
$\tfrac{1}{2}$ & & $0$ &
  Signal exactly offsets noise; zero net correction \\
$\left(-1, \tfrac{1}{2}\right)$ & Perverse & $\left(-\tfrac{3}{2}, 0\right)$ &
  Perverse beliefs; investor leans against true innovation \\
$-1$ & $-W_{t}'$ & $-\tfrac{3}{2}$ &
  Perfect misalignment; maximum growth destruction \\
\bottomrule
\end{tabular}
\end{table}

\begin{remark}[The asymmetry of intelligence]
The interval $[-\tfrac{3}{2}, +\tfrac{1}{2}]$ is not symmetric around
zero, and the asymmetry is worth sitting with. A perfectly aligned
investor ($\rho = 1$) gains $+\tfrac{1}{2}$. A perfectly
contra-indicator investor ($\rho = -1$) loses $\tfrac{3}{2}$. Bad
beliefs are three times as costly as good beliefs are valuable. The
continuous-time framework quantifies what Brennan-Lo states
qualitatively: subpopulations with $\rho < 0$ do not survive. Here
we see exactly why. The perfectly misaligned investor allocates
against every market movement, compounding her losses at each
instant. Extinction pressure, made precise.
\end{remark}

% --- BOUNDED RATIONALITY -----------------------------------------------------
\section{Bounded Rationality}

\subsection{The cost of intelligence}

In Brennan and Lo~[\hyperlink{ref:3}{3}], intelligence carries a biological cost that
caps the evolutionarily stable $\rho$. In markets, accurate beliefs
cost money: data, research, time, and attention. Let $c(\rho)$ be
the cost per unit time of achieving belief accuracy $\rho$. The net
growth rate of the dollar population is
\begin{equation}
  \bar{g}_{\mathrm{net}}(\rho) \;=\; \mub
  + \frac{1}{2}\!\left(\frac{\mue}{\sigma}\right)^{\!2}
  + \frac{\rho - \frac{1}{2}}{\dd t} - c(\rho).
  \label{eq:netgrowth}
\end{equation}
The evolutionarily stable $\rho^{*}$ maximizes
$\bar{g}_{\mathrm{net}}(\rho)$. When $c(\rho)$ is increasing and
convex, the solution sits in the interior: $\rho^{*} < 1$. Bounded
rationality is the outcome, not an assumption.

\subsection{The evolutionarily stable portfolio}

Conjecture~\ref{conj:rhomax} introduces a horizon dependence:
$\rho$ is properly written $\rho(\dd t)$, with feasibility set
$[-1, \rho_{\max}(\dd t)]$. The investor chooses $\rho(\dd t)$ to
maximize $\bar{g}_{\mathrm{net}}(\rho)$ subject to that feasibility
constraint. The first-order condition becomes
\begin{equation}
  c'\!\bigl(\rho^{*}(\dd t)\bigr) \;=\; \frac{1}{\dd t},
  \qquad
  \rho^{*}(\dd t) \;\leq\; \rho_{\max}(\dd t).
  \label{eq:foc-rho}
\end{equation}
The interior solution $c'(\rho^{*}) = 1/\dd t$ holds whenever the
unconstrained optimum lies below the feasibility ceiling. When it
does not, when the marginal benefit $1/\dd t$ would call for a $\rho$ beyond what the horizon allows, the solution is the corner
$\rho^{*}(\dd t) = \rho_{\max}(\dd t)$, and the investor is
intelligence-bound by physics rather than by cost.

Investors with low intelligence costs (quantitative funds, well-resourced institutions) reach $\rho^{*}(\dd t)$ near the upper
boundary and hold portfolios close to Kelly. Investors with high
intelligence costs (retail investors, infrequent rebalancers) land
at lower $\rho^{*}(\dd t)$ and drift from Kelly in predictable
directions.

\subsection{Coexistence of strategies}

Different investors are good at different horizons. A high-frequency
strategy and a value strategy do not compete on the same axis. Each
carries an implicit best operating point in $\dd t$. We capture this
by writing each investor's intelligence as a \emph{$\rho$-$\dd t$
curve}: a function $\rho_{j}(\dd t)$ that gives investor $j$'s
achievable correlation as a function of horizon.

The curves have characteristic shapes. The HFT investor's
$\rho_{\text{HFT}}(\dd t)$ peaks at small $\dd t$ and decays as
$\dd t$ grows. Her edge dissolves at long horizons. The value
investor's $\rho_{\text{value}}(\dd t)$ does the opposite, climbing
with $\dd t$ as fundamentals reassert themselves over noise. Each
investor's curve traces the horizons where her strategy translates
information into alignment, and the horizons where it does not.

The cross-section of investors generates a population-average curve
$\bar{\rho}(\dd t) = \E_{j}[\rho_{j}(\dd t)]$, which traces the
typical level of intelligence required to be successful at horizon
$\dd t$. The relationship between an individual investor's
$\rho_{j}(\dd t)$ and the population average $\bar{\rho}(\dd t)$ at
her chosen horizon is what determines whether she earns the
intelligence premium or pays the penalty. Investors do not compete
with one another for a single $\rho^{*}$. They compete within the
horizon they have selected, against the average investor at that
horizon.

When intelligence costs vary across investors, different investors
will optimally choose different levels of $\rho^{*}$ and will
therefore achieve different net growth rates
$\bar{g}_{\mathrm{net}}(\rho^{*})$. The portfolio of an investor who
has a lower cost of intelligence, achieving $\rho^{*}$ close to her
ceiling, will in general grow faster than the portfolio of an
investor who has a high cost of intelligence, constrained to a
lower $\rho^{*}$. The result is heterogeneous optima, each investor
growing at the rate her cost structure and her $\rho$-$\dd t$ curve
permit.

Each investor's $\rho^{*}$ is the best achievable given her cost
structure and her chosen horizon. No investor can unilaterally
improve her net growth rate by changing her strategy without also
changing her cost structure. The distribution of $\rho^{*}$ values
across investors therefore reflects the joint distribution of
intelligence costs and horizon-specific capabilities in the
population, and these distributions persist as long as their
heterogeneity persists.

This provides a measured evolutionary interpretation of the
empirical coexistence of passive indexers, active managers, and
quantitative strategies in financial markets. Each strategy reflects
a different point on the cost-benefit frontier of intelligence and a
different shape of $\rho_{j}(\dd t)$. Higher-cost strategies are
suboptimal only relative to a lower cost structure that may not be
available to the investor pursuing them. The mix of strategies
observed in the market reflects the underlying distribution of
intelligence costs and horizon-specific edges across investors,
with no single strategy positioned to drive the others to extinction.

% --- INVESTMENT HORIZON ------------------------------------------------------
\section{The Investment Horizon and the Short-Run Conundrum}

\subsection{The long-standing conundrum}

Kelly dominates in the long run~[\hyperlink{ref:1}{1}]. In the short run, Kelly can
underperform and does so with considerable volatility. The profession
has generally treated this as a fact of life. The framework here
gives it a precise resolution.

\subsection{Two distinct optimal allocations}

There are two optimal allocations, not one:
\begin{align}
  \lopt &\;=\; \frac{\mue}{\sigma^{2}}
    + \frac{1}{\sigma}\,\frac{\dd W_{t}}{\dd t},
  \label{eq:lstar2} \\[6pt]
  \lkelly &\;=\; \frac{\mue}{\sigma^{2}}.
  \label{eq:lkelly}
\end{align}
$\lopt$ maximizes instantaneous stochastic growth. $\lkelly$ maximizes
the ensemble growth rate and the long-run time-average. They are
related by
\begin{equation}
  \lim_{\dd t \to \infty} \lopt \;=\; \E[\lopt] \;=\; \lkelly.
  \label{eq:convergence}
\end{equation}
The equality is precise and carries two distinct meanings. The first
limit eliminates uncertainty by letting time run out: as $\dd t \to
\infty$, the stochastic term $(1/\sigma)\,\dd W_{t}/\dd t \sim
\mathcal{N}(0, 1/\dd t)$ vanishes, and the instantaneous
growth-optimal rule converges to $\lkelly$ pathwise. The second
equality eliminates uncertainty by aggregation: taking the expectation
over an infinite number of dollar draws, the law of large numbers
collapses the distribution to its mean, which is also $\lkelly$. Two
routes, same destination. The long-run time-average and the ensemble
average coincide, and both equal the Kelly weight, a result that
Peters~[\hyperlink{ref:11}{11}] identifies as the central insight of ergodicity economics.
The conundrum arises because the literature applies $\lkelly$ where
$\lopt$ belongs: in the short run, where the stochastic term is large
and neither limiting procedure has had a chance to operate.

\subsection{The randomization term as a hedge}

The variance of the stochastic component of $\lopt$ is
\begin{equation}
  \Var\!\left[\frac{1}{\sigma}\,\frac{\dd W_{t}}{\dd t}\right]
  \;=\; \frac{1}{\dd t \cdot \sigma^{2}}.
  \label{eq:var}
\end{equation}
The stochastic term is a hedge. Each dollar draws from a Gaussian
centered at the Kelly weight, so the expected allocation is always $\lkelly$, but the spread around that center is set by the horizon.
A shorter horizon means a wider Gaussian, a larger hedge, more
variability in how individual dollars are deployed.

\begin{proposition}[Horizon-dependent hedging]
The optimal spread of the Gaussian coin flip
\eqref{eq:gaussiancoin}, measured by its variance, is monotonically
decreasing in the investment horizon $\dd t$:
\begin{itemize}
  \item Short horizon ($\dd t$ small): variance $1/(\dd t\cdot\sigma^{2})$
    is large; each dollar hedges aggressively around Kelly; the
    dollar population fans out widely across allocation values.
  \item Long horizon ($\dd t$ large): variance shrinks; the hedge
    tightens; dollar allocations cluster near $\lkelly$.
  \item Infinite horizon ($\dd t \to \infty$): variance $\to 0$; no
    hedge required; every dollar allocates exactly at $\lkelly$.
\end{itemize}
In all cases, $\E[\lopt_{i}] = \lkelly$. The hedge does not shift
the center; it widens the spread.
\end{proposition}

Because no investor reaches infinite horizon, the hedge is always
positive. The horizon $\dd t$ is itself a preference parameter, an investor's choice governed by goals, mandate, and life stage. But
once it is fixed, the framework determines the magnitude of the
hedge as a function of $\dd t$ and $\sigma$. The decomposition is
sharp: preference enters only through the choice of horizon; the
optimal hedge for that horizon follows mechanically.

\subsection{Connection to the Brennan-Lo framework}

In Brennan-Lo, $f^{*}$ is the evolutionary imperative under
systematic risk. The term $(1/\sigma)\,\dd W_{t}/\dd t$ plays that
same role: the growth-optimal response to systematic $\dd W_{t}$,
calibrated to the investor's horizon. The Gaussian coin flip
$\mathcal{N}(\mue/\sigma^{2}, 1/(\dd t\cdot\sigma^{2}))$ is
simultaneously the evolutionary randomization mechanism of
Section~5 and the optimal hedge for any investor with a finite
horizon. Longer horizons tighten the hedge; shorter horizons widen
it. The evolutionary framework and the financial framework agree on
exactly how wide.

% --- DISCUSSION --------------------------------------------------------------
\section{Discussion}

\subsection{The normative and the evolutionary}

Kelly and Merton describe what evolutionary selection produces at the
dollar level, in a market with one risky asset and complete systematic
correlation. The investor that survives is the one who can avoid
injecting perverse intelligence into their allocation.

\subsection{The population-of-dollars reframing as a general method}

The reframing from investors to dollars may travel beyond the
single-asset problem. Many questions in financial economics involve
systematic shocks that stall the evolutionary analysis at the agent
level. Dropping to a sub-agent population (dollars, trades, individual decisions) may open those problems up. The connection to
the Brennan-Lo framework is available whenever the right population
is found.

\subsection{Implications for the Adaptive Markets Hypothesis}

The bounded rationality result in Section~8 puts numbers on the
Adaptive Markets Hypothesis~[\hyperlink{ref:8}{8}]. When excess returns $\mue$ are high,
the payoff to accurate beliefs rises, more investors pay $c(\rho)$
to push $\rho^{*}$ toward $1$, and the market grows more efficient.
When $\mue$ falls or $\sigma$ climbs, the calculus reverses,
investors settle for lower $\rho^{*}$, and efficiency erodes. The
intelligence correction $\rho - \tfrac{1}{2}$ in \eqref{eq:optlogdyn}
tracks this in real time.

\subsection{Limitations and future work}

Five extensions are immediate. First, multiple risky assets: a
vector of weights and a full covariance matrix will enrich both the
intelligence mapping and the hedging results. Second, time-varying
parameters: letting $\mue$ and $\sigma$ move connects the framework
to the return predictability literature. Third, price impact: if
investor behavior feeds back into asset prices, strategic
interactions emerge that the present framework sets aside.

Fourth, leverage constraints: the framework permits $\ell \in
\mathbb{R}$, which lets each dollar's allocation $\lopt_{i}$ take
values from $-\infty$ to $+\infty$. No real investor can use
arbitrary amounts of leverage. Future work would impose realistic
bounds on $\ell$ and examine how those bounds reshape the optimal
hedge. In particular, how the truncation of the Gaussian coin flip
distorts the convergence to the Kelly weight at finite horizons.

Fifth, the granularity of the dollar population: the framework treats
each dollar as an independent allocation decision. In practice, the
cost to enter a long or short position is almost always at least
two orders of magnitude greater than a single dollar, which makes
per-dollar allocation infeasible. Future work would examine how to
implement this strategy under a minimum-position-size constraint, in
which the effective unit of selection is not a dollar but a block of
dollars large enough to clear transaction costs. The interesting
question is how the size of that block, a function of trading costs, account size, and rebalancing frequency, governs how closely
the implementable hedge can approximate the theoretical one.

Each of these extensions is a worthwhile next step.

% --- CONCLUSION --------------------------------------------------------------
\section{Conclusion}

We extended the Brennan-Lo binary-choice framework to continuous-time
portfolio allocation with one risky and one risk-free asset. Four
results follow.

\textbf{Contribution 1.} For the individual investor, the right
population is dollars, not investors. Each dollar faces the same
systematic $\dd W_{t}$, so the Brennan-Lo logic applies directly:
randomize. The Gaussian coin flip
$\lopt_{i} \sim \mathcal{N}(\mue/\sigma^{2}, 1/(\dd t\cdot\sigma^{2}))$
is the continuous-time expression of that imperative.

\textbf{Contribution 2.} Kelly and Merton fall out of evolutionary
selection. The optimal allocation
$\lopt = \mue/\sigma^{2} + (1/\sigma)\,\dd W_{t}/\dd t$
recovers the Kelly weight in its deterministic part and encodes the
investor's beliefs about market randomness in its stochastic part.

\textbf{Contribution 3.} Intelligence is $\rho = \Corr[M_{t}',
W_{t}']$. The intelligence correction $\rho - \tfrac{1}{2} \in
[-\tfrac{3}{2}, +\tfrac{1}{2}]$ is asymmetric. Misaligned beliefs
destroy three times what aligned beliefs build. The asymmetry is
exact, and it is not gentle.

\textbf{Contribution 4.} $\lopt$ and $\lkelly = \mue/\sigma^{2}$
are different objects. The stochastic term is a hedge that each
dollar executes by drawing from a Gaussian centered at the Kelly
weight, with variance $1/(\dd t\cdot\sigma^{2})$. The expected
allocation remains Kelly regardless of horizon, but the spread around
it depends on $\dd t$. Investors with longer horizons hold tighter
distributions; investors with shorter horizons hold wider ones. The
choice of horizon belongs to the investor, but once that choice is
made, the framework specifies the hedge exactly.

\newpage
% --- REFERENCES --------------------------------------------------------------
\section{Further Reading and References}

\begin{itemize}\setlength{\itemsep}{0.4em}

  \item[{[1]}]\hypertarget{ref:1}{} Breiman, L. (1961). ``Optimal gambling systems for
    favorable games.'' \textit{Proceedings of the Fourth Berkeley
    Symposium on Mathematical Statistics and Probability}, 1, 65--78.

  \item[{[2]}]\hypertarget{ref:2}{} Brennan, T.~J. and Lo, A.~W. (2011). ``The origin of
    behavior.'' \textit{Quarterly Journal of Finance}, 1(1), 55--108.

  \item[{[3]}]\hypertarget{ref:3}{} Brennan, T.~J. and Lo, A.~W. (2012). ``An evolutionary
    model of bounded rationality and intelligence.'' \textit{PLoS ONE},
    7(11), e50310.

  \item[{[4]}]\hypertarget{ref:4}{} Cooper, W.~S. and Kaplan, R.~H. (1982). ``Adaptive
    `coin-flipping': a decision-theoretic examination of natural
    selection for random individual variation.'' \textit{Journal of
    Theoretical Biology}, 94(1), 135--151.

  \item[{[5]}]\hypertarget{ref:5}{} Farmer, J.~D. (2002). ``Market force, ecology and
    evolution.'' \textit{Industrial and Corporate Change}, 11(5),
    895--953.

  \item[{[6]}]\hypertarget{ref:6}{} Cover, T.~M. (1984). ``An algorithm for maximizing
    expected log investment return.'' \textit{IEEE Transactions on
    Information Theory}, 30(2), 369--373.

  \item[{[7]}]\hypertarget{ref:7}{} Kelly, J.~L. (1956). ``A new interpretation of
    information rate.'' \textit{Bell System Technical Journal},
    35(4), 917--926.

  \item[{[8]}]\hypertarget{ref:8}{} Lo, A.~W. (2004). ``The adaptive markets hypothesis:
    market efficiency from an evolutionary perspective.''
    \textit{Journal of Portfolio Management}, 30(5), 15--29.
    See also: Lo, A.~W. (2017). \textit{Adaptive Markets: Financial
    Evolution at the Speed of Thought}. Princeton University Press.

  \item[{[9]}]\hypertarget{ref:9}{} Merton, R.~C. (1969). ``Lifetime portfolio selection
    under uncertainty: the continuous-time case.'' \textit{Review of
    Economics and Statistics}, 51(3), 247--257.

  \item[{[10]}]\hypertarget{ref:10}{} Merton, R.~C. (1971). ``Optimum consumption and
    portfolio rules in a continuous-time model.'' \textit{Journal of
    Economic Theory}, 3(4), 373--413.

  \item[{[11]}]\hypertarget{ref:11}{} Peters, O. (2019). \textit{Ergodicity Economics}.
    London Mathematical Laboratory.

  \item[{[12]}]\hypertarget{ref:12}{} Robson, A.~J. (2001). ``The biological basis of
    economic behavior.'' \textit{Journal of Economic Literature},
    39(1), 11--33.

  \item[{[13]}]\hypertarget{ref:13}{} Robson, A.~J. (2001). ``Why would nature give
    individuals utility functions?'' \textit{Journal of Political
    Economy}, 109(4), 900--914.

  \item[{[14]}]\hypertarget{ref:14}{} Rogers, A.~R. (1994). ``Evolution of time preference
    by natural selection.'' \textit{American Economic Review},
    84(3), 460--481.

  \item[{[15]}]\hypertarget{ref:15}{} Samuelson, L. (2001). ``Introduction to the evolution
    of preferences.'' \textit{Journal of Economic Theory}, 97(2),
    225--230.

  \item[{[16]}]\hypertarget{ref:16}{} Thorp, E.~O. (1971). ``Portfolio choice and the Kelly
    criterion.'' \textit{Proceedings of the Business and Economics
    Section of the American Statistical Association}, 215--224.

\end{itemize}
\newpage
% --- DISCLOSURES -------------------------------------------------------------
\section{Disclosures}

The views expressed in this document are those of the author and
SVVVL alone and do not constitute investment, legal, tax, or
accounting advice. This document is provided for informational and
research purposes only and is not an offer or solicitation to buy or
sell any security or financial product.

The framework presented is theoretical and relies on idealized
assumptions that do not hold in practice. The results
should not be implemented in a real portfolio without additional
constraints and considerations not addressed in this paper.

While reasonable care has been taken in the preparation of this
document, no representation or warranty, express or implied, is made
as to its accuracy or completeness.

Reading this document does not create an advisory, fiduciary, or
client relationship of any kind between the reader and SVVVL or the
author.

\end{document}